\section{模型假设与符号说明}
\subsection{模型假设}

结合赛题附件给出的参数口径与山区应急救援的实际情形，本文作如下假设：

\begin{enumerate}[label=\textbf{假设 \arabic*}]
	\item \textbf{航段净空假设}：无人机在航段上的巡航海拔取该航段所经过 DEM 像元的\textbf{最高地面高程 $+50$\,m}（注意不是两端点高程的较大者），从而保证全程安全净空。
	\item \textbf{作业高度假设}：调度中心 O01 的作业高度取其地面海拔；各服务区的作业高度取地面海拔 $+30$\,m。连续访问多个服务区时，每次投送后均\textbf{从 30\,m 作业高度重新爬升}。
	\item \textbf{下降能耗假设}：附件给出的“下降能耗效率”为 0，故本文\textbf{不单独计算下降附加能耗}，下降段仅计入时间。
	\item \textbf{巡航能耗口径假设（关键）}：附件仅给出运输机的空载/满载标准航程与电池可用能量，\textbf{未给巡航功率}。本文按“飞满一个标准航程恰好耗尽一个可用能量”这一自洽口径\textbf{由航程反推巡航能耗}：$E^{hor}=(d/\Lg(q))\Euse$。该口径使返航安全约束退化为 $d\le(1-\rhog)\Lg(q)$，物理意义清晰。
	\item \textbf{爬升效率假设}：附件列名为“爬升能耗效率”（取 0.72），本文按\textbf{效率}使用，即 $E^{up}=mgh^{+}/\eta_{up}$。
	\item \textbf{返航安全假设}：每个架次完成全部投送后必须返回 O01，返航为\textbf{空载}；架次总能耗须满足 $E_{p}^{T}\le(1-\rhog)\Euse$。
	\item \textbf{货箱不可拆分假设}：货箱是整个投送与组批的最小单位，不可拆分、不可重复投送，同一货箱只能由一个架次送达一次。
	\item \textbf{能源资源假设}：共享电池与中继能源组件均为\textbf{独立资源}并各自记录 SOC，初始 SOC 为 100\%，再次投入任务前须充至 100\%；同一机型的共享电池可在该机型不同实体机间调度，\textbf{不同机型不可混用}。
	\item \textbf{通信判定假设}：链路可用性按“双向链路预算 $\wedge$ 无地形遮挡”判定，门限取两方向较小值；通信状态按“\textbf{直连 $>$ 中继 $>$ 中断}”三态优先判定，中继要求接入段与回传段\textbf{同一时刻同时可用}，且\textbf{不允许多跳转发}。
	\item \textbf{连续通信假设}：运输无人机在爬升、巡航、下降及投送四个阶段均应保持通信，因此通信判定须沿完整轨迹\textbf{逐时刻采样}，不能只查服务区一点或航段两端点。
	\item \textbf{中继共享假设}：题目附录只限制“每架\textbf{运输无人机}在任一时刻只能由 G01 或\textbf{一架}中继无人机提供通信保障”，并\textbf{未}限制一架中继的服务对象数量；据此本文允许\textbf{同一悬停点被多架运输机共享}。中继须在其保障的每个中断区间\textbf{起点之前完成建链}。
	\item \textbf{分区刚性假设}：问题四中，问题三已确定的货箱组批、服务区访问顺序、运输与中继任务安排及通信保障关系\textbf{保持不变}；同一运输架次涉及的服务区必须划入同一任务组；任务执行期间各类资源\textbf{不得跨组调配}。
\end{enumerate}

\subsection{符号说明}

\begin{table}[htbp]
	\caption{主要符号说明}
	\label{tab:symbol}
	\centering
	\zihao{-5}
	\begin{tabular}{llll}
		\toprule
		符号 & 含义 & 符号 & 含义 \\
		\midrule
		O01 & 临时调度中心 & $S_{i}$ & 第 $i$ 个服务区（$i=1,\dots,15$） \\
		G01 & 固定网关（设于 O01） & $\mathcal{G}$ & 运输机型集合（$|\mathcal{G}|=3$） \\
		$b$ & 货箱编号（$b=1,\dots,80$） & $m_{b},v_{b}$ & 货箱质量、体积 \\
		$\Qg$ & 机型 $g$ 最大载货质量 & $V_{g}$ & 机型 $g$ 最大装载体积 \\
		$L_{g0},L_{gF}$ & 空载、满载标准航程 & $\Euse$ & 机型 $g$ 单组电池可用能量 \\
		$v_{g}^{\uparrow},v_{g}^{c},v_{g}^{\downarrow}$ & 爬升、巡航、下降速度 & $\rhog$ & 返航安全余量比例 \\
		$\Tfull$ & 等效完全充电时间 & $H_{op}(n)$ & 节点 $n$ 的作业高度 \\
		$d_{ij}$ & 航段水平距离 & $h_{ij}^{+},h_{ij}^{-}$ & 航段爬升、下降高度 \\
		$t_{gij},E_{gij}(q)$ & 航段飞行时间、能耗 & $p$ & 一个运输架次 \\
		$E_{p}^{T}$ & 架次 $p$ 总运输能耗 & $q_{max}^{safe}$ & 最大安全载荷 \\
		$k_{s}$ & 第 $s$ 站投送箱数 & $t_{h0},t_{h1}$ & 基础交接时间、每箱增量时间 \\
		$d_{b}^{fb},d_{b}^{exp}$ & 首批截止时间、期望送达时间 & $t_{b}^{deliver}$ & 货箱 $b$ 实际交付时刻 \\
		$N$ & 总架次数 & $T_{\max}$ & 全部任务完成时间 \\
		$b_{ijt}$ & 地形遮挡变量（0/1） & $L_{FSPL}$ & 自由空间传播损耗 \\
		$L_{path},L_{\max}$ & 总路径损耗、最大允许损耗 & $A_{ijt}$ & 链路可用性（0/1） \\
		$P$ & 中继悬停点 & $E_{R}$ & 中继架次能耗 \\
		$K$ & 任务分区组数 & $U$ & 原子单元（连通分量） \\
		\bottomrule
	\end{tabular}
\end{table}
